\documentclass[aspectratio=169,11pt,t]{beamer}
\input{../../../../theme/ksa-cs-v2.tex}
% 원본 덱 (ksa-theme): slides/v2/ml1/week10/cnn/ -- 그림은 거기 figs/ 를 그대로 쓴다
\setbeamerfont{author}{size=\normalsize}
\setbeamerfont{institute}{size=\small}
\graphicspath{{../../../../v2/ml1/week10/cnn/figs/}{../../../../assets/}}
\makeatletter
\newcommand{\markpage}[1]{\protected@write\@auxout{}{\string\newlabel{#1}{{\the\c@framenumber}{}}}\hypertarget{#1}{}}
\newcommand{\pref}[1]{\hyperlink{#1}{\textcolor{ksagolddk}{(p.\ref{#1})}}}
\makeatother

\renewcommand{\ksacredit}{Machine Learning and Deep Learning \,\textperiodcentered\, CNN 합성곱 신경망}
\title{CNN 합성곱 신경망 (Convolutional Neural Network)}
\subtitle{작은 필터 하나를 그림 전체에 재사용한다}
\author{Machine Learning and Deep Learning}
\institute{한국과학영재학교 (KSA)}
\date{}
\begin{document}
{\setbeamertemplate{footline}{}\begin{frame}[c,noframenumbering]\titlepage\end{frame}}

\section{질문}
% --- 오늘의 질문
\begin{frame}{오늘의 질문}
\vfill
\centering
{\LARGE 손글씨 숫자 하나 (28 × 28 픽셀) 를 알아보는 데}\\[0.6em]
{\LARGE 가중치 \kq{125만 개}가 꼭 필요할까?}\\[1.2em]
{\large 완전연결: 픽셀 784개를 뉴런 1,600개에 전부 연결 \quad $784 \times 1{,}600 + 1{,}600 = 1{,}256{,}000$}
\vfill
\note{지금까지의 신경망은 모든 입력을 모든 뉴런에 연결하는 완전연결 신경망이었습니다. 28 곱하기 28 손글씨 그림은 픽셀이 784개입니다. 첫 층의 뉴런이 1600개면 가중치는 784 곱하기 1600에 편향 1600을 더해 125만 6천 개입니다. 숫자 하나 알아보는 데 정말 이렇게 많이 필요할까요? 오늘은 훨씬 적은 가중치로 그림을 더 잘 보는 신경망, 합성곱 신경망, CNN을 배웁니다.}
\end{frame}

\section{합성곱}
% --- 완전연결의 문제
\begin{frame}{완전연결의 두 가지 문제}
\vfill
\centering
\renewcommand{\arraystretch}{1.6}
{\large\begin{tabular}{ll}
\textbf{가중치가 너무 많다} & 그림이 커지면 폭발 (224 × 224 컬러면 입력만 15만 개)\\
\textbf{위치를 모른다} & 같은 가로획이 한 칸 옮겨지면 \kq{완전히 다른 입력}으로 본다\\
 & 이웃한 픽셀끼리 가깝다는 사실을 쓰지 않는다
\end{tabular}}\vskip0.8em
{\large 그림의 성질: \kb{가까운 픽셀끼리 의미가 있다} \quad \kb{같은 모양은 어디에 있든 같은 모양}}
\vfill
\note{완전연결에는 두 가지 문제가 있습니다. 첫째, 가중치가 너무 많습니다. 손글씨는 작아서 그나마 괜찮지만 224 곱하기 224 컬러 사진이면 입력만 15만 개라서 가중치가 폭발합니다. 둘째, 그림의 구조를 쓰지 않습니다. 완전연결은 픽셀 784개를 그냥 숫자 784개로 보기 때문에, 같은 가로획이 한 칸 오른쪽으로 옮겨지면 완전히 다른 입력으로 봅니다. 하지만 그림에는 두 가지 성질이 있습니다. 가까운 픽셀끼리 모여 의미를 만들고, 같은 모양은 어디에 있든 같은 모양입니다. CNN은 이 두 성질을 그대로 쓴 구조입니다.}
\end{frame}

% --- 합성곱
\begin{frame}{합성곱: 3 × 3 필터를 겹쳐 곱하고 더한다}
\markpage{conv}
\vfill
\centering\includegraphics[width=0.84\linewidth]{cn_window.png}\par\vskip0.1em
{\large $(1 + 1 + 1) \times (+1) + (0.4 + 0.3 + 0.4) \times 0 + (0 + 0 + 0) \times (-1) = \kb{3.0}$}\\[0.2em]
{\small 위가 밝고 아래가 어두우면 큰 값 = 가로선의 윗변 \quad 필터 = 숫자 9개}
\vfill
\note{CNN의 핵심은 필터입니다. 숫자 9개짜리 3 곱하기 3 필터를 그림의 한 부분에 겹치고, 겹친 칸끼리 곱해서 모두 더합니다. 이것을 합성곱이라고 부릅니다. 예로 위쪽은 플러스 1, 가운데는 0, 아래쪽은 마이너스 1인 가로선 필터를 7의 가로획 위, 빨간 창에 겹쳐 봅니다. 위 줄 픽셀은 모두 1이라 3, 가운데 줄은 0을 곱해 0, 아래 줄은 픽셀이 0이라 0, 더하면 3.0입니다. 위가 밝고 아래가 어두운 곳, 즉 가로선의 윗변에서 큰 값이 나옵니다. 이 필터는 설명을 위해 사람이 고른 것이고, 실제 CNN은 필터 값을 학습으로 정합니다.}
\end{frame}

% --- 특징 지도
\begin{frame}{필터를 한 칸씩 옮기며 그림 전체를 훑는다}
\vfill
\centering\includegraphics[width=0.8\linewidth]{cn_maps.png}\par\vskip0.2em
{\large \textbf{특징 지도} (feature map): 가로선 필터는 가로획이, 세로선 필터는 \kb{내려오는 획}이 밝다}\\[0.2em]
{\small 어느 위치에서든 \textbf{같은 9개}를 쓴다 = 파라미터 공유 \quad 같은 모양은 어디에 있든 같은 반응}
\vfill
\note{같은 필터를 한 칸씩 옮기며 그림 전체를 훑으면 결과도 그림 모양이 됩니다. 이것을 특징 지도라고 합니다. 가로선 필터의 특징 지도는 7의 가로획이 밝고, 세로선 필터의 특징 지도는 비스듬히 내려오는 획이 밝습니다. 중요한 점은 어느 위치에서든 같은 숫자 9개를 쓴다는 겁니다. 이를 파라미터 공유라고 합니다. 그래서 가로획이 위에 있든 아래에 있든 같은 필터가 똑같이 반응합니다. 완전연결의 두 문제, 가중치 수와 위치를 한꺼번에 해결한 겁니다.}
\end{frame}

\section{크기 · 가중치}
% --- 출력 크기
\begin{frame}{출력 크기}
\markpage{size}
\vfill
\centering
{\LARGE\[ \text{출력} = \left\lfloor \frac{n - k + 2p}{s} \right\rfloor + 1 \]}\vskip0.4em
\renewcommand{\arraystretch}{1.35}
{\large\begin{tabular}{llc}
$n = 28$, $k = 3$, 패딩 $p = 0$, 보폭 $s = 1$ & 가장자리 한 칸씩 줄어든다 & \kb{26}\\
$p = 1$ (가장자리에 0 을 한 겹) & 크기 유지 (``same'') & 28\\
$p = 1$, $s = 2$ (두 칸씩 건너뛰기) & 절반으로 & 14
\end{tabular}}\vskip0.4em
{\small $k$: 필터 크기 \quad $p$: 패딩 (padding) \quad $s$: 보폭 (stride)}
\vfill
\note{출력 크기는 이 식으로 구합니다. 그림 크기 n에서 필터 크기 k를 빼고 패딩의 두 배를 더해 보폭으로 나눈 뒤 1을 더합니다. 28 곱하기 28 그림에 3 곱하기 3 필터를 패딩 없이 한 칸씩 옮기면 28 빼기 3 더하기 1, 26입니다. 3 곱하기 3 창이 그림 밖으로 나갈 수 없어서 가장자리가 한 칸씩 줄어듭니다. 가장자리에 0을 한 겹 둘러싸는 패딩 1을 쓰면 크기가 28로 유지됩니다. 보폭을 2로 해서 두 칸씩 건너뛰면 14로 절반이 됩니다.}
\end{frame}

% --- 채널
\begin{frame}{채널과 가중치 수}
\vfill
\centering
{\large 입력 채널 $c_{\text{in}}$ 개 (컬러면 3) \quad $\to$ 필터 하나 = $k \times k \times c_{\text{in}}$ \quad 필터 $c_{\text{out}}$ 개 = 출력 채널}\vskip0.4em
{\LARGE\[ \text{가중치} = (k^2\,c_{\text{in}} + 1)\,c_{\text{out}} \]}\vskip0.2em
\renewcommand{\arraystretch}{1.3}
{\large\begin{tabular}{lcr}
1층: 1 $\to$ 4 채널 & $(9 \cdot 1 + 1) \times 4$ & 40\\
2층: 4 $\to$ 8 & $(9 \cdot 4 + 1) \times 8$ & 296\\
3층: 8 $\to$ 16 & $(9 \cdot 8 + 1) \times 16$ & 1,168\\\hline
합계 & & \kb{1,504}
\end{tabular}}
\vfill
\note{실제 CNN에는 필터가 여러 개 있습니다. 필터 하나가 특징 지도 하나를 만들고, 필터 수만큼의 특징 지도를 출력 채널이라고 부릅니다. 다음 층은 이 여러 채널을 입력으로 받으니, 다음 층 필터는 3 곱하기 3 곱하기 입력 채널 수의 상자 모양입니다. 컬러 사진은 처음부터 빨강, 초록, 파랑 3채널입니다. 가중치 수는 필터 하나의 숫자 k 제곱 곱하기 입력 채널에 편향 1을 더하고 필터 수를 곱하면 됩니다. 채널을 4, 8, 16으로 늘리는 세 층은 40, 296, 1168, 합쳐서 1504개입니다.}
\end{frame}

% --- 비교
\begin{frame}{1,504 개 vs 1,256,000 개}
\vfill
\centering
{\Large CNN 3층 \kb{1,504개} \quad vs \quad 완전연결 \kq{1,256,000개} \quad (835배)}\\[0.8em]
\includegraphics[width=0.62\linewidth]{cn_compare.png}\par\vskip0.2em
{\large 실제 학습: 작은 CNN \kb{5,258개로 97.5\%} \quad 완전연결 101,770개로 96.4\%}
\vfill
\note{비교해 봅시다. 채널 4, 8, 16인 CNN 세 층은 가중치 1504개입니다. 같은 크기의 결과를 완전연결로 만들면 125만 6천 개, 약 835배입니다. 실제로 학습해 보면 어떨까요? MNIST에서 합성곱 층 두 개짜리 작은 CNN은 가중치 5258개로 2에폭 만에 97.5퍼센트, 완전연결 신경망은 약 20배 많은 10만 개로 96.4퍼센트입니다. 그림의 구조를 쓰니 적은 가중치로도 더 잘 봅니다.}
\end{frame}

\section{풀링 · 구조}
% --- 풀링
\begin{frame}{풀링: 작게 줄이면서 특징은 남긴다}
\markpage{pool}
\vfill
\centering\includegraphics[width=0.56\linewidth]{cn_pool.png}\par\vskip0.2em
{\large \textbf{최대 풀링}: 2 × 2 중 \kb{가장 큰 값}만 \quad 26 × 26 $\to$ 13 × 13 \quad 가중치 0개}\\[0.2em]
{\small 크기는 $\frac14$, 계산량도 줄어든다 \quad 획이 한 칸쯤 움직여도 결과가 비슷하다}
\vfill
\note{CNN에는 풀링이라는 층도 있습니다. 가장 흔한 최대 풀링은 2 곱하기 2 네 칸 중 가장 큰 값만 남깁니다. 26 곱하기 26 특징 지도가 13 곱하기 13이 되어 크기는 4분의 1이지만, 가로획은 그대로 살아 있습니다. 풀링에는 학습할 가중치가 없고, 다음 층의 계산량을 줄여 줍니다. 또 네 칸 중 최댓값만 보니 획이 한 칸쯤 움직여도 결과가 거의 같아서, 조금 비뚤게 쓴 숫자에도 강해집니다.}
\end{frame}

% --- 구조
\begin{frame}{전형적인 CNN 구조}
\vfill
\centering
{\Large 입력 $\to$ [합성곱 $\to$ ReLU $\to$ 풀링] $\times$ 몇 번 $\to$ 펼치기 $\to$ 완전연결 $\to$ 출력}\\[1.0em]
\renewcommand{\arraystretch}{1.4}
{\large\begin{tabular}{ll}
앞쪽 층 & 선 · 모서리 같은 작은 특징\\
뒤쪽 층 & 앞 특징의 조합: 고리, 획의 교차, 눈 · 바퀴 같은 부분\\
마지막 & 완전연결로 ``이 특징들이면 7''
\end{tabular}}\vskip0.6em
{\small LeNet-5 (LeCun 외, 1998): 우편번호 숫자 인식 \quad AlexNet (2012): 이미지넷 1위, 깊은 CNN 시대의 시작}
\vfill
\note{CNN은 보통 합성곱, ReLU, 풀링을 묶어 몇 번 반복한 뒤, 마지막에 특징 지도를 한 줄로 펼쳐 완전연결층으로 답을 냅니다. 앞쪽 층은 선이나 모서리 같은 작은 특징을, 뒤쪽 층은 그 조합인 고리나 교차, 더 깊으면 눈이나 바퀴 같은 부분을 찾습니다. 풀링으로 크기가 줄어들수록 뒤쪽 필터 하나가 원래 그림의 더 넓은 영역을 보게 되기 때문입니다. 1998년 르쿤의 LeNet이 이 구조로 우편번호 숫자를 읽었고, 2012년 AlexNet이 이미지넷 대회에서 1위를 하며 깊은 CNN 시대가 열렸습니다.}
\end{frame}

% --- 학습된 필터
\begin{frame}{학습된 첫 층 필터}
\vfill
\centering\includegraphics[width=0.9\linewidth]{cn_filters.png}\par\vskip0.4em
{\large MNIST 로 학습한 작은 CNN 의 첫 층 필터 8개 (빨강 +, 파랑 $-$)}\\[0.4em]
{\large 사람이 정해 주지 않았는데 \kb{가로 · 세로 · 대각선 모서리}를 찾는 모양이 된다}
\vfill
\note{실제로 학습한 CNN의 첫 층 필터 8개를 그려 봤습니다. 빨강이 양수, 파랑이 음수입니다. 앞에서 사람이 고른 가로선 필터처럼, 위아래나 좌우, 대각선으로 밝고 어두운 경계를 찾는 모양들이 보입니다. 아무도 가로선을 찾으라고 알려 주지 않았는데, 숫자를 맞히는 데 필요한 특징을 스스로 배운 겁니다.}
\end{frame}

\section{연습}
% --- 문제풀이
\begin{frame}{문제풀이}
\vfill
\begin{enumerate}
\item 32 × 32 그림에 5 × 5 필터, 패딩 0, 보폭 1 이면 출력 크기는? 패딩 2 면?\vskip0.8em
\item 입력 3채널 (컬러), 3 × 3 필터 64개인 합성곱 층의 가중치 수는?\vskip0.8em
\item 픽셀 값 $\begin{bmatrix}0 & 0 & 0\\ 1 & 1 & 1\\ 1 & 1 & 1\end{bmatrix}$ 에 가로선 필터 (위 +1, 가운데 0, 아래 $-1$) 를 쓰면?\vskip0.8em
\item 특징 지도 $\begin{bmatrix}1 & 3\\ 2 & 0\end{bmatrix}$ 에 2 × 2 최대 풀링을 하면?
\end{enumerate}
\vfill
\note{1번 답. 32 빼기 5 더하기 1, 28입니다. 패딩 2면 32 그대로입니다. 2번 답. 9 곱하기 3 더하기 1, 28을 64개 곱해 1792개입니다. 3번 답. 위 줄 0, 아래 줄 3에 마이너스 1을 곱해 마이너스 3입니다. 아래가 밝은 가로선의 아랫변이라 음수가 나옵니다. 4번 답. 3입니다.}
\end{frame}

% --- 실습
\begin{frame}{실습}
\vfill
\centering\begin{tabular}{rl}
\large 1 & 합성곱 · 최대 풀링 직접 구현 $\to$ 손글씨 7 에서 3.0, 26 × 26, 13 × 13 재현\\[0.5em]
\large 2 & 출력 크기 · 가중치 수 계산 함수 $\to$ PyTorch 와 비교\\[0.5em]
\large 3 & 작은 CNN vs 완전연결 (MNIST): 가중치 수와 정확도, 학습된 필터 그리기
\end{tabular}\vskip1em
MNIST (LeCun, Cortes, Burges; torchvision) \\[0.4em]\href{https://colab.research.google.com/github/SuminHan/book-ml/blob/main/slides/v2/ml1/week10/cnn/lab/cnn.ipynb}{\textcolor{ksagolddk}{▶ Colab에서 실습 노트북 열기}}
\vfill
\note{실습은 세 부분입니다. 첫째, 합성곱과 최대 풀링을 numpy로 직접 구현해서 손글씨 7에서 3.0, 26 곱하기 26, 13 곱하기 13을 재현합니다. 둘째, 출력 크기와 가중치 수를 계산하는 함수를 만들어 PyTorch 층과 비교합니다. 셋째, 작은 CNN과 완전연결 신경망을 MNIST로 학습해서 가중치 수와 정확도를 비교하고, 학습된 첫 층 필터를 그려 봅니다.}
\end{frame}

\section{정리}
% --- 정리
\begin{frame}{정리}
\vfill
\centering\large\begin{tabular}{l}
\textbf{합성곱}: 작은 필터를 겹쳐 곱하고 더하기, 그림 전체를 훑어 특징 지도\\[0.5em]
\textbf{파라미터 공유}: 어디서든 같은 필터 $\to$ 1,504 vs 1,256,000 (835배)\\[0.5em]
출력 크기 $\frac{n - k + 2p}{s} + 1$ \quad 가중치 $(k^2 c_{\text{in}} + 1) c_{\text{out}}$\\[0.5em]
\textbf{최대 풀링}: 2 × 2 중 최댓값, 크기 $\frac14$ \quad 필터는 학습으로 모서리를 찾는다
\end{tabular}
\vfill
\note{정리합니다. 합성곱은 작은 필터를 그림에 겹쳐 곱하고 더하는 계산이고, 그림 전체를 훑으면 특징 지도가 나옵니다. 어디서든 같은 필터를 쓰는 파라미터 공유 덕분에 같은 크기의 완전연결보다 835배 적은 가중치로 그림을 봅니다. 출력 크기와 가중치 수는 오늘의 두 식으로 계산합니다. 최대 풀링은 2 곱하기 2 중 최댓값만 남겨 크기를 4분의 1로 줄입니다. 그리고 필터 값은 학습으로 정해지며, 스스로 모서리를 찾는 모양이 됩니다.}
\end{frame}
\end{document}
